\section{Common rank tests on the selected history}
\label{sec:actual-rank-purchases}

We now organize the comparison estimates into the finite set of tests
needed to prevent loss of rank. The essential quantifier is that one
calibration serves every row of the testing matrix. Its choice may depend
on the comparison scale and on the finite volume, but cannot depend on
the row subsequently tested. We also retain the original partition
function throughout the auxiliary integration history.

Work with the funded calibrated family for the simple gauge group of the
construction, retaining its paid calendar and local comparison inputs.
Fix a grade $d$. Let $V_d$ be its reference coefficient space, with its
real structure, distinguished unit $e_0$, and coefficient norm $|\cdot|_1$.
Fix a finite real basis $q_0,\ldots,q_r$ of the free quotient, starting
with the scalar, and real representatives $v_0,\ldots,v_r\in V_d$,
with $v_0=e_0$. Such representatives exist: lift each real quotient vector
and replace a lift $v$ by $(v+\bar v)/2$. The quotient projection commutes
with conjugation, so its image is unchanged. The scalar has its prescribed
lift. Fix also the finite generic measurement configurations separating
these basis vectors. Each row uses either one or two insertions.

The argument below uses the same-family local estimates of the preceding
construction: the corrected mixed compiler, its reader comparisons, the
calibrated insertion prices, and the constant-coupling auxiliary history.
Their numerical inequalities are imposed at the selected comparison
scales. In particular the assertion that the compiler applies to every
matched palette and that the actual reader has its required mean and
covariance bounds remains part of this input. The finite-dimensional
selection in this section does not supply a missing local estimate.

\subsection{Buying comparison scales before selecting the history}

Write $j$ for a stage of the construction and $k$ for a level of the
realized calendar; $J(k)$ is its stage. The preliminary purchase supplies
a finite menu $I_j$ of grades at stage $j$, with every fixed grade belonging
to $I_j$ for all sufficiently large $j$. For the finite requests in $I_j$
choose the comparison parameters and then the necessary block floor.
Two kinds of bounds enter this floor: those ensuring the local estimates
and those allowing their comparison parameters to be bought. Their
maximum, denoted $\mu_j$, serves both purposes. Increasing the preliminary
block prescription preserves these properties. Choose $L_j\ge\mu_j$.

Some parameters can only be chosen after $L_j$. For a grade $i\in I_j$,
let $a_{ij}>0$ be the admissible upper bound for its final marker radius.
Let $c_{ij}>0$ be the radius already allowed by its comparison data.
After fixing a common positive outer parameter for the finite menu, choose
\[
             0<\delta_{ij}\le\min(a_{ij},c_{ij},1).
\]
The local comparison is stable under this reduction of the marker radius.
Choose its remaining threshold after $\delta_{ij}$ and $L_j$. Taking the
maximum over $i\in I_j$ gives a single eventual coupling threshold for
the stage. No maximum over infinitely many grades is required.

The paid calendar is enlarged at these thresholds and at the existing
window and minimum-length bounds. The realization property of that
calendar says that every sufficiently late stage occurs on a paid run
of levels, with its block, coupling and window data. This proves that,
for each fixed grade, suitable comparison levels are unbounded. Indeed,
given a lower level $k_0$, choose a stage $j>J(k_0)$ in whose menu the
grade occurs, also beyond the finitely many construction floors. Realize
that stage at a level $k$. If $k\le k_0$, monotonicity of $J$ would give
$j=J(k)\le J(k_0)$, a contradiction. Thus $k>k_0$.

At these levels all the local comparisons use the same price map and
the same geometric history. To check this last point, the preliminary
shell purchase is enlarged to
\[
                     \text{shell}_j\ge100(j+7)L_j.
\]
The radius needed by the comparison is $(j+7)L_j$, and the paid volume
bound therefore supplies its factor-$100$ room. The path of the selected
history is the path of the purchased evolution at every finite cutoff.
Its funding and support-link records consequently give the comparison
on that path with precisely the original price map. This is the point
at which a local estimate is attached to the history: a second
independently realized paid calendar would not justify the substitution.

In the selected construction, these choices occur in the following
order. The finite menus and block floors are chosen first; the enlarged
shell delays and post-length thresholds come next; the common paid
calendar is realized next. Finally the callback which receives the
selected finite-volume history constructs its comparison binding using
the path identity just checked. Thus subsequent use of a frequent set
of comparison levels does not choose another trajectory.

\subsection{One readout and one auxiliary history for all rows}

Retain only comparison levels $k\ge\operatorname{ref}(d)$ and discard
the finitely many levels where the purchased exponent $I_k$ below is
negative. This is possible since $I_k\to\infty$. Fix such a suitable
level $k$ and write $s$ for its clock coupling. The actual
couplings obey the purchased bound
\[
                  |\beta_{N,p,k}-s|\le s/8
                  \qquad(N\ge k).
\]
Here $N$ is the mesh cutoff and $p$ the period index. At every such finite
volume the construction of Section~\ref{sec:restarted-rank-history}
gives a constant-coupling auxiliary history on the required window.
Its transition identities and all its local comparisons refer to the
same endpoint as the original history. Its additive remainder need only
have the value bounds stated there.

The calibrated insertion bounds hold eventually in $N$, uniformly for
every $p$, and for every row in the finite testing family. Intersect these
finitely many eventual conditions before choosing the auxiliary history.
The local readout construction then supplies a single complex-linear map
$T_{k,N,p}:V_d\to V_d$. Write $A_{k,N,p}$ for the primary-coordinate map,
including its factor $a_k^{-d}$, where $a_k>0$ is the lattice spacing.
The readout satisfies
\begin{equation}\label{eq:rank-common-tested}
 T_{k,N,p}e_0=e_0,\qquad
 T_{k,N,p}\bar v=\overline{T_{k,N,p}v},\qquad
 |A_{k,N,p}v-T_{k,N,p}v|_1\le\tfrac12|A_{k,N,p}v|_1.
\end{equation}
These are estimates on the reference coefficient space. They are obtained
with the same reader and mixed coefficient estimates used for the marked
integrals, not by choosing an arbitrary near-invertible map afterwards.

Here is how the previously proved comparison steps combine at this
choice. The common coordinate construction of
Section~\ref{sec:allarity-coordinate-calibration} gives the readout and
its covariance atlas. The actual covariance estimates of
Section~\ref{sec:actual-covariance-transport} apply to this atlas at the
prescribed configurations. The paid coefficient bound for every named
lift makes its terminal insertion admissible. For every row, the
healthy-field argument of Section~\ref{sec:marked-healthy-cumulant-launch}
then compares the actual middle-scale polynomial to the free polynomial.
Section~\ref{sec:allarity-cumulant-comparison} performs the connected
coefficient subtraction. The same-row marked propagation of
Section~\ref{sec:allarity-marked-propagation} transports this comparison
through the lower part of the auxiliary window. All finite row and
column requests are imposed before the eventual threshold is chosen.
The comparison radius and local constants can depend on $d$ and on this
finite family, but the subsequent threshold does not depend on $p$.

In more precise terms, this produces the following joint statement at
unboundedly many $k$. For all sufficiently large $N$, and every $p$, there
are one map $T_{k,N,p}$ satisfying \eqref{eq:rank-common-tested}, one
integration course, and one auxiliary representation at the lower end
of the window. For every family of lifts solving the normalization
equations below, the coefficient estimates hold for every row and column
on this same representation. The marked density identity keeps its
partition function equal to that of the original endpoint. The
quantifier over normalized lifts follows the choice of the readout;
there is no separate readout for different lift families.

Only arities at most two enter this assertion. At the middle scale the
local healthy comparison starts the required one- and two-marker
coefficients. The propagation keeps their invoice, meaning the original
marked density after subtraction of its prescribed free connected
coefficients. The normalization estimate in
Section~\ref{sec:rank-root-normalization} applies at the lower end of
that very integration course. If $I\ne\varnothing$ is a chosen marker
subset and $q$ is the number of terminal cells, its estimate is
\[
 a^{q|I|}|\mathbb E(\text{replicated coefficient})|
       \le \exp(q I_k),\qquad a=\exp(2I_k).
\]
The finite-volume harmonic correction is already included in $I_k$ as
specified there. Taking the $q$th root gives
\begin{equation}\label{eq:rank-common-root}
       \|\text{coefficient }I\|_{\rm root}
       \le e^{I_k-2I_k|I|}\le e^{-I_k}=:\epsilon_k.
\end{equation}
The clock purchase makes $I_k\to\infty$, so $\epsilon_k\to0$.
The empty coefficient is the unaltered scalar normalization and is not
subject to this estimate. In the two-marker coefficient the singleton
subtractions are retained; replacing it by an unsubtracted two-point
observable would not give \eqref{eq:rank-common-root}.

\subsection{Solving the normalization equations exactly}

The reference transport estimates give a number $D_k>0$, independent of
$N\ge k$ and $p$, such that
\begin{equation}\label{eq:rank-reference-inverse}
                       |v|_1\le D_k|A_{k,N,p}v|_1.
\end{equation}
Explicitly one may take
\[
 D_k=a_k^d\prod_{u=0}^{k-1} C_{u+1}L_u^d,
\]
where $C_{u+1}\ge1$ is the purchased front-column bound. This follows
by applying that bound successively from the normalized reference level
to $k$; the extra earlier factors are at least one. The factor $a_k^d$
undoes the normalization in $A_{k,N,p}$.

For brevity suppress $k,N,p$ in this paragraph. By the triangle inequality
and \eqref{eq:rank-common-tested},
\[
                         |Av|_1\le2|Tv|_1.
\]
If $Tv=0$, \eqref{eq:rank-reference-inverse} implies $v=0$.
Thus $T$ is injective and hence bijective on the finite-dimensional space
$V_d$. This also proves the inverse estimate $|v|_1\le2D_k|Tv|_1$.

Each coordinate of $V_d$ has two nonnegative integer degrees: a position
degree and a factor degree. Denote multiplication by $\rho$ to the first
degree by $W_\rho$, and multiplication by $t$ to the second by $S_t$.
They are distinct diagonal maps. Put
\[
 t_{k,N,p}=\frac{\beta_{N,p,k}^{-1/2}}{\theta_k}>0,
 \qquad \rho_k=\mathcal M_k^b,
\]
where $\theta_k$ is the original clock threshold, $\mathcal M_k$ its
monitor and $b>0$ the fixed comparison exponent. For every $i$ solve
\begin{equation}\label{eq:rank-common-lifts}
                  S_{t_{k,N,p}}T_{k,N,p}u_i=W_{\rho_k}v_i.
\end{equation}
There is exactly one solution, since both maps on the left are
invertible. Because all the maps commute with conjugation and $v_i$ is
real, conjugating the equation and using uniqueness shows $\bar u_i=u_i$.
Both diagonal maps fix $e_0$, so uniqueness and $Te_0=e_0$ give $u_0=e_0$.
All columns are therefore normalized simultaneously, with their actual
unit and real structure.

The same equations give a bound independent of the two cutoffs at fixed
$k$. The coupling band and the clock identity
$\theta_k=\widehat\beta_k^{-1/2}\mathcal M_k$ imply
\[
            t_{k,N,p}^{-1}
            =\theta_k\sqrt{\beta_{N,p,k}}\le2\mathcal M_k.
\]
Since all factor degrees are nonnegative,
\[
 |T_{k,N,p}u_i|_1
 =|S_{t_{k,N,p}^{-1}}W_{\rho_k}v_i|_1
 \le |S_{2\mathcal M_k}W_{\rho_k}v_i|_1.
\]
Together with the inverse estimate this bounds every $|u_i|_1$ by
\begin{equation}\label{eq:rank-fixed-scale-bounded}
            B_k=2D_k\sum_{i=0}^r
                        |S_{2\mathcal M_k}W_{\rho_k}v_i|_1.
\end{equation}
No uniform bound in $k$ is asserted or needed. For the finitely many mesh
cutoffs preceding the common eventual threshold, define every lift to be
$e_0$. Enlarging $B_k$ by $|e_0|_1$ makes the bound valid for all $N,p$,
without changing any eventual marked estimate.

\subsection{Selecting a sequence and taking its ordered limits}

Let $\ell(k)$ denote the lower end of the comparison window. On the
admissible comparison levels the calendar is constant in stage throughout
$[\ell(k),k]$. Its depth tends to infinity, and its physical test scale is
\[
                       \sigma_k=a_k\mathcal M_k^b\longrightarrow0.
\]
The generic lattice choices for each fixed measurement configuration
converge, after division by $\sigma_k$, to that configuration. Their
physical coordinates lie in the middle half of a cell at level $\ell(k)$;
in particular they have the fixed relative boundary margin $1/4$.
These are the geometric fit and tuning conditions purchased for the
finite configurations. They hold eventually at admissible levels.

Intersect the unbounded set of jointly funded comparison levels with
these eventual conditions and with the eventual normalization conditions
of Section~\ref{sec:rank-root-normalization}. For any $k_0,\ell_0$ and
$\eta>0$ there remains a level $k$ with
\[
                  k>k_0,\qquad \ell(k)>\ell_0,
                  \qquad \sigma_k<\eta.
\]
For the middle inequality, choose a stage later than $J(\ell_0)$.
The plateau property forces its lower window end past $\ell_0$.
The other two inequalities follow from cofinality and $\sigma_k\to0$.
Choose $k_0$ in this set, and recursively choose $k_{n+1}$ with
\[
 k_{n+1}>k_n,\qquad \ell(k_{n+1})>\ell(k_n),
                      \qquad \sigma_{k_{n+1}}<\sigma_{k_n}.
\]
The physical scales are strictly decreasing and tend to zero, because
$k_n\to\infty$. The error bounds $\epsilon_{k_n}$ also tend to zero;
their monotonicity is unnecessary. Choose the simultaneous normalized
lifts just constructed at each $k_n$. Their cut-off uniform bounds are
those of \eqref{eq:rank-fixed-scale-bounded}.

At each $n$ and all sufficiently large $N$, uniformly for every period,
the insertion stencils lie in the comparison cell and every nonempty
corrected coefficient has the root bound \eqref{eq:rank-common-root}.
There are only finitely many rows and columns at fixed $d$, so their
thresholds have a common maximum. The resulting sequence thus supplies
the bounded lifts, geometry and moment estimates needed for the common
extraction of Section~\ref{sec:finite-rank-kernel-transport}.

Include all these bounded lift sequences, for every grade and comparison
index, in that same countable extraction. First take the mesh limit for
each fixed period, then the period limit. Denote the limiting lifts by
$u_i^{(n)}$. Reality and the unit identity survive these limits. The
finite-cell moment inequality and the corrected root bound give, for
every row and column,
\begin{equation}\label{eq:rank-purchased-point-comparison}
 \left|S(\sigma_{k_n}x^{(n)};,u^{(n)})
                 -S_{\rm free}(x^{(n)};,q)\right|
             \le P_n\epsilon_{k_n}.
\end{equation}
Here $S$ is the extracted interacting correlation, $q$ lists the free
basis elements used by the row, and $x^{(n)}$ is its rescaled lattice
configuration. The factor $P_n$ is the finite polynomial in the absolute
free subset moments appearing when the connected subtractions are
undone. It is bounded: $x^{(n)}$ tends to the fixed generic configuration,
all its subsets remain separated there, and the free correlations are
continuous on those domains. Hence the right side tends to zero.

To justify the order of this passage, fix $n$ first. Its root bound holds
for every period and all mesh cutoffs past one threshold. Apply the
finite-cell inequality at those cutoffs; joint convergence of the lifts
and correlations gives the inequality after the mesh limit at each
period, and then after the period limit. Only then let $n\to\infty$.
There is no interchange between the comparison-scale limit and either
finite-volume limit. Equation~\eqref{eq:rank-purchased-point-comparison}
is exactly the finite measurement convergence used in the rank argument.
Identification of the entire equation kernel additionally requires the
all-spectator local equation comparison stated in that argument; finite
one- and two-point measurements alone do not supply it.
